\section{Illustrative site investigation}
\subsection{Field programme}
The assumed investigation includes three rotary wash boreholes (BH-01 to BH-03), standard
penetration testing at approximately \SI{1.5}{\metre} intervals, disturbed and thin-wall samples,
and temporary standpipes. For a real project, the final scope should be set after reviewing
column locations, basement geometry and access constraints.

\begin{table}[H]
\centering
\caption{Illustrative exploration schedule.}
\label{tab:exploration}
\begin{tabularx}{\textwidth}{@{}L{20mm}C{24mm}C{25mm}C{31mm}Y@{}}
\toprule
\rowcolor{navy!8}
\textbf{Point} & \textbf{Depth} & \textbf{Ground level} & \textbf{Water level} & \textbf{Principal tests} \\
 & (m) & (m datum) & (m below grade) & \\
\midrule
BH-01 & 25.0 & 100.20 & 2.1 & SPT, classification, strength \\
BH-02 & 25.0 & 100.05 & 2.3 & SPT, consolidation, chemistry \\
BH-03 & 25.0 &  99.90 & 2.2 & SPT, classification, strength \\
\bottomrule
\end{tabularx}
\end{table}

\subsection{Subsurface profile}
\begin{figure}[H]
\centering
\resizebox{\textwidth}{!}{%
\begin{tikzpicture}[x=1cm,y=0.36cm,font=\footnotesize]
  % depth axis
  \draw[-{Stealth[length=2mm]},thick,navy] (0,1) -- (0,-26);
  \node[rotate=90,navy] at (-0.75,-12.5) {Depth below grade (m)};
  \foreach \d in {0,5,10,15,20,25}{
    \draw[navy] (-0.12,-\d)--(0.12,-\d);
    \node[left] at (-0.15,-\d) {\d};
    \draw[muted!25,dashed] (0.2,-\d)--(15.8,-\d);
  }
  % main profile
  \fill[fillsoil] (1,0) rectangle (5,-1.2);
  \fill[claysoil] (1,-1.2) rectangle (5,-4.5);
  \fill[sandsoil] (1,-4.5) rectangle (5,-10.5);
  \fill[densesand] (1,-10.5) rectangle (5,-18.0);
  \fill[claysoil!75!navy] (1,-18.0) rectangle (5,-25.0);
  \draw[navy,thick] (1,0) rectangle (5,-25);
  \foreach \y in {-1.2,-4.5,-10.5,-18}{\draw[white,line width=1pt] (1,\y)--(5,\y);}
  \node[align=center] at (3,-0.6) {FILL};
  \node[align=center] at (3,-2.85) {SOFT--FIRM\\SILTY CLAY};
  \node[align=center] at (3,-7.5) {MEDIUM-DENSE\\SILTY SAND};
  \node[align=center] at (3,-14.25) {DENSE SAND};
  \node[align=center,text=white] at (3,-21.5) {STIFF CLAY\\WITH SAND};
  % water line
  \draw[blue,densely dashed,line width=1.2pt] (0.5,-2.2)--(15.8,-2.2);
  \node[blue,fill=white,inner sep=1pt] at (13.9,-2.2) {GWT $\approx$ \SI{2.2}{m}};
  % pile schematic
  \fill[navy!80] (7.3,-2.7) rectangle (7.85,-17.5);
  \fill[navy] (7.05,-2.35) rectangle (8.10,-2.7);
  \draw[gold,line width=2pt,{Stealth[length=2mm]}-{Stealth[length=2mm]}]
    (8.65,-2.7)--node[right,align=left]{\SI{600}{mm} bored pile\\$L=\SI{14.8}{m}$ below cap}(8.65,-17.5);
  \draw[teal,dashed] (5.1,-17.5)--(10,-17.5);
  \node[teal,anchor=west] at (9.1,-17.5) {preliminary toe};
  % notes
  \node[anchor=west,align=left,text width=3.4cm] at (12,-9) {
    \textbf{Design implication}\\[2pt]
    Soft clay controls shallow-foundation settlement.\\[4pt]
    Dense sand below \SI{10.5}{m} provides the primary pile bearing stratum.\\[4pt]
    Maintain at least \SI{5}{m} penetration into dense sand.};
\end{tikzpicture}
}
\caption{Generalized ground model and recommended pile concept. Layer boundaries may vary
between exploration points.}
\label{fig:ground-model}
\end{figure}

\subsection{Groundwater}
The interpreted groundwater level is \SIrange{2.1}{2.3}{\metre} below grade. A single reading
after drilling may be affected by the drilling method and does not define the seasonal high.
Piezometers should therefore be installed and monitored before basement and temporary-works
design. For preliminary uplift and construction checks, use a conservative groundwater level of
\SI{1.0}{\metre} below grade unless monitoring demonstrates a lower design level.

% =============================================================================
